\section{Программа <<Предки>>}

\subsection{База знаний}
Исходный код базы знаний <<Предки>> приведен в листинге~\ref{lst:t1-prog}.
\begin{listing}[Исходный код базы знаний <<Предки>>]{lst:t1-prog}{Prolog}
find_grandparent(GP, GPGender, PGender, Child) :- 
    parent(GP, Parent),
    gender(GP, GPGender),
    parent(Parent, Child),
    gender(Parent, PGender).

parent(ivan, anna).
parent(ivan, boris).
parent(maria, anna).
parent(maria, boris).

parent(petr, dmitry).
parent(petr, elena).
parent(olga, dmitry).
parent(olga, elena).

parent(anna, kirill).
parent(dmitry, kirill).
parent(boris, svetlana).
parent(elena, svetlana).

gender(ivan, male).
gender(petr, male).
gender(boris, male).
gender(dmitry, male).
gender(kirill, male).

gender(maria, female).
gender(olga, female).
gender(anna, female).
gender(elena, female).
gender(svetlana, female).
\end{listing}

Исходный код базы знаний <<Максимум>> приведен в листинге~\ref{lst:t2-prog}.
\begin{listing}[Исходный код базы знаний <<Максимум>>]{lst:t2-prog}{Prolog}
max3_cut(X, Y, Z, X) :- X >= Y, X >= Z, !.
max3_cut(_, Y, Z, Y) :- Y >= Z, !.
max3_cut(_, _, Z, Z).

max3(X, Y, Z, X) :- X >= Y, X >= Z.
max3(X, Y, Z, Y) :- Y > X,  Y >= Z.
max3(X, Y, Z, Z) :- Z > X,  Z > Y.

max2(X, Y, X) :- X >= Y.
max2_(X, Y, Y) :- X < Y.

max2_cut(X, Y, X) :- X >= Y, !.
max2_cut(_, Y, Y).
\end{listing}

\subsection{Вопросы задания}

\subsubsection{База знаний «Предки»}

Для получения имён всех бабушек заданного субъекта необходимо 
задать вопрос: существует ли такие имена предков, удовлетворяющие 
отношению \code{find\_grandparent}, при условии, что пол предка 2-го колена
конкретизирован как женский, имя внука конкретизировано.

\begin{listing}[Поиск всех бабушек]{lst:q1}{Prolog}
?- find\_grandparent(GP, female, \_, ivan).
\end{listing}

Для получения имён всех дедушек заданного субъекта необходимо 
задать вопрос: существует ли такие имена предков, удовлетворяющие 
отношению \code{find\_grandparent}, при условии, что пол предка 2-го колена
конкретизирован как мужской, имя внука конкретизировано.

\begin{listing}[Поиск всех дедушек]{lst:q2}{Prolog}
?- find\_grandparent(GP, male, \_, ivan).
\end{listing}

Для получения имён всех предков 2-го колена заданного субъекта необходимо 
задать вопрос: существует ли такие имена предков, удовлетворяющие 
отношению \code{find\_grandparent}, при условии, что имя внука конкретизировано.

\begin{listing}[Поиск всех предков 2-го колена]{lst:q3}{Prolog}
?- find\_grandparent(GP, \_, \_, ivan).
\end{listing}

Для получения имён всех бабушек по материнской линии заданного субъекта необходимо 
задать вопрос: существует ли такие имена предков, удовлетворяющие 
отношению \code{find\_grandparent}, при условии, что пол предка 2-го колена
конкретизирован как женский, пол предка 1-го колена
конкретизирован как женский, имя внука конкретизировано.

\begin{listing}[Бабушка по материнской линии]{lst:q4}{Prolog}
?- find\_grandparent(GP, female, female, ivan).
\end{listing}

Для получения имён всех бабушек по материнской линии заданного субъекта необходимо 
задать вопрос: существует ли такие имена предков, удовлетворяющие 
отношению \code{find\_grandparent}, при условии, что пол предка 1-го колена
конкретизирован как женский, имя внука конкретизировано.

\begin{listing}[Предки по материнской линии]{lst:q5}{Prolog}
?- find\_grandparent(GP, \_, female, ivan).
\end{listing}

\subsubsection{Поиск максимума}

Для нахождения максимального из двух чисел 
необходимо задать вопрос о существовании такой переменной \code{Max}, 
которая удовлетворяет отношению \code{max2} или \code{max2\_cut} при 
конкретизированных первом и втором числе.

\begin{listing}[Максимум из двух чисел]{lst:q6}{Prolog}
?- max2\_cut(10, 20, Max).
\end{listing}

Для нахождения максимального из трех чисел 
необходимо задать вопрос о существовании такой переменной \code{Max}, 
которая удовлетворяет отношению \code{max3} или \code{max3\_cut} при 
конкретизированных первом, втором и третьем числе.

\begin{listing}[Максимум из трех чисел]{lst:q7}{Prolog}
?- max3\_cut(5, 12, 7, Max).
\end{listing}